% Generated draft: requires mathematical and source review.
\documentclass[11pt]{article}
\usepackage[margin=25mm]{geometry}
\usepackage{fontspec}
\setmainfont{Latin Modern Roman}
\usepackage{amsmath,amssymb,mathtools,unicode-math}
\setmathfont{Latin Modern Math}
\usepackage{enumitem,needspace,xurl,hyperref}
\setlength{\parindent}{0pt}
\setlength{\parskip}{5pt}
\setlength{\emergencystretch}{3em}
\newcommand{\N}{\mathbb{N}}
\newcommand{\Z}{\mathbb{Z}}
\newcommand{\Q}{\mathbb{Q}}
\newcommand{\R}{\mathbb{R}}
\newcommand{\C}{\mathbb{C}}
\newcommand{\F}{\mathbb{F}}
\newcommand{\A}{\mathbb{A}}
\newcommand{\PP}{\mathbb{P}}
\newcommand{\E}{\mathbb{E}}
\newcommand{\eps}{\varepsilon}
\newcommand{\dd}{\,\mathrm{d}}
\begin{document}
\section*{UnsolvedMath ID: 2583}
\textbf{Source problem code:} catalogue entry
\subsection{Definitions and mathematical statement}
The source asks three separate decision questions for an input $g\in\operatorname{CT}(\Z)$. It does not define the group $\operatorname{CT}(\Z)$ or specify a finite effective representation of its elements. These data are essential to an algorithmic decidability statement and are not reconstructed from the notation or from the review's informal phrase ``compatible transformations''.

Conditionally regarding $g$ as a permutation of $\Z$, a residue class is a set $a+m\Z$ with $a\in\Z$ and integer $m\geq1$. To permute a partition $\mathcal P$ means $\{g(C):C\in\mathcal P\}=\mathcal P$, with distinct parts disjoint and union $\Z$. Task (a) asks whether such an invariant nontrivial residue-class partition exists. The source does not specify whether the partition must be finite, whether its parts can have different moduli, or the exact nontriviality convention. In particular, no restriction to all classes of one fixed modulus is silently inserted.

For $z\in\Z$, its cycle is its orbit under the permutation $g$. The two cycle predicates are precise once the action is known:
\begin{itemize}
\item (b) Only finite cycles: $\forall z\in\Z\ \exists k\in\Z_{\geq1}$ with $g^k(z)=z$. There is no required uniform bound on $k$, so this is not the same as asking that $g$ have finite order.
\item (c) No finite cycles: $\forall z\in\Z\ \forall k\in\Z_{\geq1}$, $g^k(z)\ne z$.
\end{itemize}
These are distinct properties, not logical complements: a permutation may have both finite and infinite cycles. For each task, decidability would mean a procedure that halts on every valid encoded group element and returns the correct yes/no answer. The admissible encodings and group action must be supplied before that claim has a fixed computational meaning.

The problem is attributed to S. Kohl in Kourovka Notebook issue 21 (2026), Problem 21.74. The supplied triage checked 2026-08-17 labels all three tasks open (OPEN-TRIAGE), says no deciding algorithms were verified, and supplies no theorem beyond recording the separate tasks. The registry is also open. No decision algorithm, undecidability result, or exact interpretation of the undefined group is invented here.
\subsection{Short English statement}
Given an effectively represented element of the source group $\operatorname{CT}(\Z)$, decide whether it preserves a nontrivial residue-class partition, has only finite cycles, or has no finite cycles. The record omits the group definition, input encoding, and partition conventions needed for a fully precise algorithmic statement.
\subsection{Sources}
\begin{itemize}
\item Local UnsolvedMath registry, permanent ID 2583, source code catalogue entry, ``Kourovka Notebook Problem 21.74''; supplied three decision tasks and attribution to S. Kohl.
\item \emph{Kourovka Notebook}, 21st issue (2026), Problem 21.74, supplied maintained source: \url{https://www.kourovka-notebook.org/}. The record supplies the homepage rather than a linked group definition or a precise input presentation.
\item Supplied literature triage checked 2026-08-17, OPEN-TRIAGE; it reports no verified deciding algorithm for any task, records no source alteration, and cautions that dated triage does not exclude later or unindexed results.
\end{itemize}
\end{document}
